\section{The ordinary double tilt throughout the real sector}
\label{sec:heart}

For the factor $\Theta=1$, we now identify the heart furnished by
Proposition~\ref{prop:geometric-prescribed-charges}.  The method of comparing
hearts through smooth surface sections follows Cheng--Feyzbakhsh
\cite[Section~3]{ChengFeyzbakhsh2026}.  Here the charge has already been
corrected to the ordinary Chern character, and positive weighted sections
allow the polarization to be real.  The argument below proves the phase
boundaries directly.

Fix $B\in V,H\in W,t\ge T$ from the preceding construction.  Let
$\mathcal P$ be the induced slicing on $X$, with charge $Z_{B,tH}$, and put
$\mathcal C=\mathcal P(0,1]$.  As before,
$H=\sum_h\ell_hL_h$ with every $\ell_h>0$.  All restrictions of complexes
in this section are derived restrictions.

\subsection{The geometric heart on a surface}

\begin{lemma}\label{lem:surface-geometric-heart}
Let $S$ be a smooth projective surface and $(Z,\mathcal R)$ a locally finite
geometric stability condition, normalized so that point sheaves have phase
$1$.  Suppose
\[
 \Im Z(G)=\omega\bigl(c_1(G)-\beta\rk(G)\bigr),
\]
where $\omega$ is a real ample divisor class and $\beta$ a real divisor
class.  Then
\[
 \mathcal R(0,1]=\langle\mathcal F_{\omega,\beta}[1],
                              \mathcal T_{\omega,\beta}\rangle,
\]
where $\mathcal T_{\omega,\beta}$ has strictly positive slope factors
(including torsion) and $\mathcal F_{\omega,\beta}$ has nonpositive
slope factors.
\end{lemma}

\begin{proof}
Lemma~\ref{lem:geometric-range} gives
$\mathcal R(0,1]\subset D^{[-1,0]}(S)$ and
$\Coh(S)\subset\mathcal R(-1,1]$.  The comparison of these two bounded
$t$-structures makes $\mathcal R(0,1]$ a tilt of $\Coh(S)$, with torsion
pair
\[
 \mathcal T_0=\Coh(S)\cap\mathcal R(0,1],\qquad
 \mathcal F_0=\Coh(S)\cap\mathcal R(-1,0].
\]
For instance their torsion sequence is obtained by truncating a sheaf at
phase $0$; the cohomological range ensures that both pieces are sheaves.

If $G\in\mathcal T_0$, every quotient sheaf also belongs to
$\mathcal T_0$ and has nonnegative imaginary charge.  A nonzero
positive-rank quotient cannot have imaginary charge zero: it would lie in
$\mathcal R(1)$, whereas Lemma~\ref{lem:geometric-range} makes its zeroth
cohomology zero-dimensional.  Thus all positive-rank quotients of $G$ have
strictly positive slope.  Applying this to its last slope
Harder--Narasimhan quotient proves
$\mathcal T_0\subset\mathcal T_{\omega,\beta}$.
If $F\in\mathcal F_0$, then $F[1]\in\mathcal R(0,1]$, so $F$ is
torsion-free by the same lemma.  Every subsheaf remains in $\mathcal F_0$
and has nonpositive imaginary charge.  Its maximal-slope factor therefore
has slope at most zero, giving
$\mathcal F_0\subset\mathcal F_{\omega,\beta}$.

These inclusions identify the torsion pairs.  Indeed the torsion sequence
for either pair, applied to an object in the corresponding part of the
other, leaves a quotient or subobject in the intersection of its two
orthogonal parts, hence zero.  This proof includes slope zero; it does not
use a generic-parameter argument.
\end{proof}

Take a smooth section $j_h:S_h\hookrightarrow X$ in $|L_h|$ and its
compatible slicing $\mathcal P_h$.  Grothendieck--Riemann--Roch for this
divisor embedding, whose normal line bundle is $L_h|_{S_h}$, gives
\begin{equation}\label{eq:surface-induced-charge}
 Z_h(G)=-\int_{S_h}e^{-B|_{S_h}-itH|_{S_h}}\ch(G)
             \frac{1-e^{-L_h|_{S_h}}}{L_h|_{S_h}}.
\end{equation}
Since $(1-e^{-L})/L=1-L/2+L^2/6$ on a surface, its imaginary part is
\begin{equation}\label{eq:surface-imaginary-charge}
 \Im Z_h(G)=tH|_{S_h}\bigl(c_1(G)
             -(B|_{S_h}+L_h|_{S_h}/2)\rk(G)\bigr).
\end{equation}
Lemma~\ref{lem:surface-geometric-heart} identifies
\begin{equation}\label{eq:section-surface-heart}
 \mathcal P_h(0,1]=\langle\mathcal F_h[1],\mathcal T_h\rangle,
\end{equation}
where the surface slope uses $H|_{S_h}$ and twist
$B|_{S_h}+L_h|_{S_h}/2$.  Only the imaginary part of the induced charge
is needed for this identification; its real part need not be the
ordinary surface charge.

\subsection{Weighted sections give strict slope bounds}

The section-phase comparison of
\cite[Lemma~3.1]{ChengFeyzbakhsh2026} applies using the line-twist
estimate \eqref{eq:deformed-twist-distance}. The two section triangles give, for every nonzero restricted object,
\begin{align}
 \phi^-_{\mathcal P_h}(E|_{S_h})
       &\ge\phi^-_{\mathcal P}(E),\label{eq:section-lower-phase}\\
 \phi^+_{\mathcal P_h}(E(L_h)|_{S_h})
       &\le\phi^+_{\mathcal P}(E)+1.\label{eq:section-upper-phase}
\end{align}
To verify the first inequality, push the triangle
$E(-L_h)\to E\to j_{h,*}(E|_{S_h})$ to $Y$.
Its last term is an extension of $E$ and $E(-L_h)[1]$, and the latter
has smallest phase at least that of $E$ by the strict twist-distance bound.
For the second, $j_{h,*}(E(L_h)|_{S_h})$ is an extension of $E(L_h)$ and
$E[1]$, both of largest phase at most $\phi^+_{\mathcal P}(E)+1$.
Pushforward computes the induced phases, proving the assertions.

For a positive-rank sheaf, use the unnormalized slope
\[
 m_{H,B}(F)=\frac{H^2(c_1(F)-B\rk F)}{\rk F};
\]
its signs agree with the normalized slopes defining
$\mathcal B_{H,B}$.  Superscripts $+$ and $-$ denote the largest and
smallest Harder--Narasimhan slopes, with torsion assigned slope $+\infty$.
The next estimate adapts the argument of
\cite[Proposition~3.3]{ChengFeyzbakhsh2026}
by combining the sections with positive real weights. Put
\begin{equation}\label{eq:weighted-slope-gap}
 c=\frac12\sum_h\ell_hHL_h^2>0.
\end{equation}

\begin{lemma}\label{lem:heart-cohomology-slopes}
Every $E\in\mathcal C$ has cohomology in degrees $[-2,0]$, its sheaf
$\mathcal H^{-2}(E)$ is torsion-free, and
\begin{equation}\label{eq:heart-sign-bounds}
 m^+_{H,B}(\mathcal H^{-2}(E))\le-c,
 \qquad m^-_{H,B}(\mathcal H^0(E))>c.
\end{equation}
A bound for a zero sheaf is vacuous.
\end{lemma}

\begin{proof}
The range and torsion-freeness follow from
Lemma~\ref{lem:geometric-range}.  Write $F=\mathcal H^{-2}(E)$ and
$G=\mathcal H^0(E)$.  For each $h$, choose a general smooth $S_h$ which
is a nonzerodivisor on the finitely many cohomology sheaves and the
subobjects and quotients used below.  Those sheaves and sequences then
restrict without Tor.  Such choices are available inside the complete
very ample linear system.

By \eqref{eq:section-upper-phase}, $E(L_h)|_{S_h}$ has largest phase at
most $2$.  The surface tilt \eqref{eq:section-surface-heart}, applied to
its lowest ordinary cohomology, implies
$F(L_h)|_{S_h}\in\mathcal F_h$.  Similarly
\eqref{eq:section-lower-phase} places $E|_{S_h}$ in
$\mathcal P_h(>0)$, and the highest ordinary cohomology of the surface
tilt gives $G|_{S_h}\in\mathcal T_h$.  These conclusions can also be
read directly from the aisles: $\mathcal P_h(\le2)$ has its degree $-2$
cohomology in $\mathcal F_h$, while $\mathcal P_h(>0)$ has its degree
$0$ cohomology in $\mathcal T_h$.

If $F\ne0$, restrict its maximal-slope subsheaf $F'\subset F$.
The class $\mathcal F_h$ is closed under subsheaves, so the shifted twist
in the definition of the surface slope gives
\[
 \frac{HL_h(c_1(F')-B\rk F')}{\rk F'}
       \le-\frac12HL_h^2.
\]
If $G$ has a positive-rank last slope quotient $G''$, quotient-closure of
$\mathcal T_h$ gives
\[
 \frac{HL_h(c_1(G'')-B\rk G'')}{\rk G''}
       >\frac12HL_h^2.
\]
For a torsion sheaf $G$ its minimum slope is $+\infty$ instead.
Multiply the displayed inequalities by the positive $\ell_h$ and sum.
Since $\sum_h\ell_hL_h=H$, they are exactly
\eqref{eq:heart-sign-bounds}.  Positivity of the weights is what permits
this argument for a real ample $H$.
\end{proof}

\subsection{Comparison with the first tilt}

Write $\mathcal B=\mathcal B_{H,B}$ for the ordinary first tilt.  The slope
bounds just obtained place it between two adjacent translates of the
constructed heart.

\begin{lemma}\label{lem:first-tilt-phase-range}
We have
\begin{equation}\label{eq:first-tilt-phase-range}
 \mathcal B\subset\mathcal P(-1,1].
\end{equation}
Consequently
\[
 \mathcal T_0=\mathcal B\cap\mathcal P(0,1],\qquad
 \mathcal F_0=\mathcal B\cap\mathcal P(-1,0]
\]
is a torsion pair in $\mathcal B$, and
$\mathcal C=\langle\mathcal F_0[1],\mathcal T_0\rangle$.
\end{lemma}

\begin{proof}
Let $D\in\mathcal T_{H,B}$ be a sheaf with $m^-_{H,B}(D)>0$.
Lemma~\ref{lem:geometric-range} places all its phases in $(-2,1]$.
If a last Harder--Narasimhan factor $G$ had phase in $(-2,-1]$, then
$G[2]\in\mathcal P(0,1]$.  It follows that $G\in D^{[0,2]}$ and
$m^+_{H,B}(\mathcal H^0(G))\le-c$.  The canonical nonzero map $D\to G$
is tested on its degree-zero cohomology:
\[
 \Hom(D,G)=\Hom(D,\mathcal H^0(G))=0.
\]
Here the equality follows from $G\in D^{\ge0}$ and standard truncation;
the vanishing is the slope comparison, including the case of torsion $D$
since the target is torsion-free.  This contradiction excludes that phase
interval and gives $D\in\mathcal P(-1,1]$.

Now let $J\in\mathcal F_{H,B}$, so $J$ is torsion-free with
$m^+_{H,B}(J)\le0$.  The coherent-sheaf range gives
$J[1]\in\mathcal P(-1,2]$.  If a first Harder--Narasimhan factor $G$
had phase in $(1,2]$, then $G[-1]\in\mathcal P(0,1]$,
$G\in D^{[-3,-1]}$, and
$m^-_{H,B}(\mathcal H^{-1}(G))>c$.
The canonical nonzero map $G\to J[1]$ is tested on its top cohomology:
\[
 \Hom(G,J[1])=\Hom(\mathcal H^{-1}(G),J)=0,
\]
again by truncation and slope.  Thus $J[1]\in\mathcal P(-1,1]$.
Extensions prove \eqref{eq:first-tilt-phase-range}.

The final assertion is the comparison theorem for two bounded hearts one
of which lies in two adjacent degrees of the other.  Explicitly, truncate
an object of $\mathcal B$ at phase $0$ in $\mathcal P$.
The inclusion \eqref{eq:first-tilt-phase-range} and its corresponding
inclusions of aisles make both truncations objects of $\mathcal B$;
they give its torsion sequence with the two stated intersections.
Tilting that sequence recovers $\mathcal P(0,1]$.
\end{proof}

\subsection{The real-axis boundary and the second tilt}

The constructed heart is now a tilt of $\mathcal B$.  It remains to
identify the torsion pair, including the distinction between positive
and zero second slope.  The following observation, following the boundary argument of
\cite[Lemma~3.6]{ChengFeyzbakhsh2026}, controls precisely those objects.

\begin{lemma}\label{lem:phase-one-first-tilt}
An object of $\mathcal B\cap\mathcal P(1)$ is a zero-dimensional sheaf.
\end{lemma}

\begin{proof}
Let $E$ belong to the stated intersection.  Its ordinary cohomology is
$J=\mathcal H^{-1}(E)\in\mathcal F_{H,B}$ and
$U=\mathcal H^0(E)\in\mathcal T_{H,B}$.
Lemma~\ref{lem:geometric-range} makes $U$ zero-dimensional.
Choose the smooth sections $S_h$ to avoid its finite support and to
restrict $J$ without Tor.  If $J\ne0$, it is torsion-free of positive
rank, and $E|_{S_h}=J|_{S_h}[1]$.
Inequality \eqref{eq:section-lower-phase} gives
$\phi^-_{\mathcal P_h}(J|_{S_h})\ge0$.
Its phases as a coherent sheaf are also at most $1$, by
Lemma~\ref{lem:geometric-range}.  Hence
$\Im Z_h(J|_{S_h})\ge0$ and \eqref{eq:surface-imaginary-charge} gives
\[
 \frac{HL_h(c_1(J)-B\rk J)}{\rk J}\ge\frac12HL_h^2
 \quad\text{for each }h.
\]
The positive weighted sum says $m_{H,B}(J)\ge c>0$, contrary to
$m^+_{H,B}(J)\le0$.  Therefore $J=0$ and $E=U$.
\end{proof}

We use the physical second slope from the introduction, writing
$\nu^{\mathrm{phys}}_{B,tH}=\nu_{B,tH}$ to distinguish it from
$\nu_{a,b}$:
\begin{equation}\label{eq:physical-second-slope}
 \nu^{\mathrm{phys}}_{B,tH}(E)=
 \frac{tH\ch_2^B(E)-t^3H^3\ch_0(E)/6}
 {t^2H^2\ch_1^B(E)},
\end{equation}
with value $+\infty$ at a zero denominator.  The tilt
Harder--Narasimhan property for each fixed real ample class $tH$ and real
$B$ is the input of \cite[Section~2]{BMS2016}; here no continuity as the
ample direction varies is required.  Its Harder--Narasimhan filtration
defines the torsion pair $(\mathcal T',\mathcal F')$ in
$\mathcal B$ with strictly positive slopes in $\mathcal T'$ and
nonpositive slopes in $\mathcal F'$.  Multiplication by a positive
constant does not affect this pair.  The claimed double tilt is
$\mathcal A_{B,tH}=\langle\mathcal F'[1],\mathcal T'\rangle$.

\begin{lemma}\label{lem:physical-zero-denominator}
For $E\in\mathcal B$ the denominator in
\eqref{eq:physical-second-slope} is nonnegative.  If it vanishes, then
$\Im Z_{B,tH}(E)\ge0$, with equality only for a zero-dimensional sheaf.
\end{lemma}

\begin{proof}
Let $J=\mathcal H^{-1}(E)$ and $U=\mathcal H^0(E)$.
The first torsion pair makes each contribution to
$H^2\ch_1^B(E)=H^2\ch_1^B(U)-H^2\ch_1^B(J)$ nonnegative.
If the sum is zero, $J$, when nonzero, is slope-semistable of slope zero,
and $U$ has dimension at most one.  Indeed positive-rank quotients in the
positive torsion part have strictly positive degree, and a nonzero
effective divisorial cycle has positive intersection with $H^2$.
For real $H$, the K\"ahler Bogomolov--Gieseker inequality follows from
\cite[Theorem~1.1, Equation~(1.6)]{LiZhangZhang2017}, applied with zero
Higgs field to the reflexive hull $J^{**}$ when $J\ne0$.
This hull is semistable for the same $H$: intersecting a subsheaf of
$J^{**}$ with $J$ preserves its rank and first Chern class.
The quotient $Q=J^{**}/J$ has codimension at least two, and, writing
$r=\rk J$, additivity gives
\[
 c_1(J)^2-2r\ch_2(J)
 =c_1(J^{**})^2-2r\ch_2(J^{**})+2r\ch_2(Q),
\]
where $\ch_2(Q)$ is an effective curve cycle, so the inequality descends
to $J$.
The same argument on a surface uses an effective zero-cycle correction;
neither application changes the real polarization.
Classical Bogomolov--Gieseker and the real-ample Hodge index theorem
\cite[Theorem~6.4]{GrebRossToma2016}, applied to
$H^2\ch_1^B(J)=0$, give
$H\ch_2^B(J)\le0$.  The curve cycle of $U$ is effective.  Consequently
\[
 \Im Z_{B,tH}(E)=
 tH\ch_2^B(U)-tH\ch_2^B(J)+\frac{t^3H^3}{6}\rk J\ge0.
\]
If $J\ne0$, the last term is strictly positive.  If $J=0$ and $U$ has a
nonzero one-dimensional part, its curve degree is strictly positive.
Thus equality is possible exactly when $E=U$ is zero-dimensional.
\end{proof}

\begin{proposition}\label{prop:exact-double-tilt-heart}
For the factor $\Theta=1$ throughout the real sector of
Proposition~\ref{prop:geometric-prescribed-charges},
\[
 \mathcal P(0,1]=\mathcal A_{B,tH}.
\]
\end{proposition}

\begin{proof}
Compare $(\mathcal T',\mathcal F')$ with the torsion pair
$(\mathcal T_0,\mathcal F_0)$ of
Lemma~\ref{lem:first-tilt-phase-range}.
For $D\in\mathcal T'$, take its latter torsion sequence
$0\to D_1\to D\to D_2\to0$.
A nonzero $D_2\in\mathcal F_0$ has $\Im Z(D_2)\le0$.
Its denominator cannot vanish: Lemma~\ref{lem:physical-zero-denominator}
would make it zero-dimensional, whereas zero-dimensional sheaves are in
$\mathcal P(1)$ and $\mathcal F_0\subset\mathcal P(-1,0]$.
Thus $D_2$ has positive denominator and
$\nu^{\mathrm{phys}}_{B,tH}(D_2)\le0$, contradicting the strictly positive
slopes of every nonzero quotient of $D$.  Hence
$\mathcal T'\subset\mathcal T_0$.

For $D\in\mathcal F'$, consider $D_1$ in the same torsion sequence.
A nonzero subobject of $D$ has second slope at most zero, so $D_1$ has
positive denominator and $\Im Z(D_1)\le0$.
But $D_1\in\mathcal T_0\subset\mathcal P(0,1]$ makes its imaginary
charge nonnegative.  It must therefore vanish, and $D_1\in\mathcal P(1)$.
Lemma~\ref{lem:phase-one-first-tilt} makes it zero-dimensional, whose second
slope is $+\infty$, a contradiction.  Thus
$\mathcal F'\subset\mathcal F_0$.
The two inclusions identify the torsion pairs, and their tilts coincide.
\end{proof}

Proposition~\ref{prop:exact-double-tilt-heart}, together with
Proposition~\ref{prop:geometric-prescribed-charges}, proves the ordinary
charge assertion of Theorem~\ref{thm:geometric-sector}, including its exact
heart, full numerical support, local finiteness, stable points, and one
threshold for an open real sector.  The equality of hearts was proved only
for $\Theta=1$.  The Todd-charge construction retains its geometric heart
from the preceding section.

Finally, the coefficient in the strong inequality should be kept distinct
from the weaker sign consequence of this construction.  If $E\ne0$ is
first-tilt semistable of finite physical slope zero, then
$E[1]\in\mathcal A_{B,tH}$ and its nonzero real charge is negative.  Thus
\[
 \ch_3^B(E)<\frac{t^2}{2}H^2\ch_1^B(E).
\]
With $t=\sqrt3a$, the strong inequality has coefficient $t^2/18$ instead.
The categorical argument supplies the real-sector and full-support
conclusions; the sharper coefficient will be proved independently in
Sections~\ref{sec:apex} and~\ref{sec:wall} by the fixed-volume estimate
and wall transport.
